% Lösungen Einheit 2 von 8 – Laplace-Experimente in der Oberstufe (zusammenbau v0.1)
\einheitenkopf{Einheit 2 von 8 · Laplace-Experimente in der Oberstufe}

\erg{14}{a) ja, Laplace \quad b) $\frac{8}{50} = \frac{4}{25}$ \quad c) Alle Murmeln sind gleich groß und gleich schwer und werden blind gezogen – keine ist bevorzugt, jede hat dieselbe Chance. Weiteres Beispiel: der Wurf einer fairen Münze. \quad d) $\frac{3}{8}$ (KKZ, KZK, ZKK) \quad e) 4 günstige von 36: (3; 6), (4; 5), (5; 4), (6; 3); $P = \frac{4}{36} = \frac{1}{9}$ \quad f) 4 und 4 günstige Ergebnisse: Summe 5 (1; 4), (4; 1), (2; 3), (3; 2), Summe 9 (3; 6), (6; 3), (4; 5), (5; 4) – gleich wahrscheinlich \quad g) 9; 4; 23 (Gewinn; Freilos; Niete) von 36 Sektoren \quad h) $P = 1$: Fall rot – 5 rote in X, 5 blaue in Y; Fall blau – 4 rote in X, 4 blaue in Y \quad i) $\frac{0{,}2n + 2}{n + 2} - 0{,}2 = \frac{1{,}6}{n + 2} > 0$ für jedes n}
\erg{15}{a) Nicht haltbar: Es fehlt die Zahl aller Räder je Stadt; die Wahrscheinlichkeit ist der Anteil der gestohlenen an allen Rädern.}
\erg{16}{a) Jan zählt die 6 Paare mit gleicher Augenzahl mit; „kleiner als“ schließt sie aus: $21 - 6 = 15$, $P = \frac{15}{36} = \frac{5}{12}$ \quad b) Beide Wahrscheinlichkeiten haben denselben Nenner, die Zahl aller Ergebnisse; größer ist die mit mehr günstigen Ergebnissen. \quad c) 2; 8; 14 (Hauptgewinn; Trostpreis; Niete) \quad d) 1; 1; 2 Kugeln (rot; blau; gelb)}
